% Szkic po polsku (konferencja do ustalenia): symulatory wyżarzania a D-Wave Advantage. Wersja anonimowa.
\documentclass[11pt]{article}
\usepackage[T1]{fontenc}
\usepackage[utf8]{inputenc}
\usepackage[polish]{babel}
\usepackage[margin=2.5cm]{geometry}
\usepackage{amsmath,amssymb}
\usepackage{booktabs}
\usepackage{graphicx}
\usepackage{url}
\usepackage[hidelinks,pdfauthor={},pdftitle={}]{hyperref}
\newcommand{\tofill}[1]{\textbf{[#1]}}

\title{Czy symulator może zastąpić wyżarzacz kwantowy? Symulacja klasyczna,
kwantowe Monte Carlo i dokładna dynamika wobec D-Wave Advantage
w przeszukiwaniu sąsiedztwa}
\author{Anonim}
\date{}

\begin{document}
\maketitle

\begin{abstract}
\tofill{Szkic: streszczenie powstanie po wynikach.}
\end{abstract}

\section{Wstęp}

Wyżarzacze kwantowe rozwiązują problemy kwadratowej optymalizacji binarnej
bez ograniczeń (QUBO), powoli wyłączając pole poprzeczne~\cite{kadowaki1998,albash2018aqc},
a komercyjne urządzenia z tysiącami kubitów są
dostępne~\cite{johnson2011manufactured}. Dostęp do nich jest limitowany, więc
w czasie prac badacze często zastępują urządzenie symulatorem: klasycznym
symulowanym wyżarzaniem, symulowanym wyżarzaniem kwantowym metodą Monte Carlo
po całkach po trajektoriach albo, dla bardzo małych układów, dokładną
dynamiką. Wcześniejsze prace porównywały takie symulatory z urządzeniami
D-Wave na losowych szkłach spinowych i pytały, czy urządzenie zachowuje się
jak wyżarzacz kwantowy, czy klasyczny~\cite{boixo2014evidence,ronnow2014speedup}.
Praktyk zadaje węższe pytanie: jeśli symulator zastąpi urządzenie wewnątrz
algorytmu optymalizacji, czy algorytm dojdzie do tego samego wyniku?

Odpowiadamy na nie dla konkretnego zadania: przeszukiwania sąsiedztwa
w permutacyjnym problemie przepływowym, w którym wyżarzacz kwantowy wybiera
ruch iteracyjnego przeszukiwania lokalnego, rozwiązując okienkowe
QUBO~\cite{wojewodzki2026windowed,wojewodzki2026bpasts}. Duży zbiór przebiegów
na D-Wave Advantage już istnieje, więc porównanie nie wymaga nowego czasu
urządzenia. Te same QUBO uruchamiamy na kilku symulatorach w trzech
scenariuszach o rosnącym realizmie, porównujemy osiągane długości
uszeregowania z wynikami urządzenia i podajemy koszt każdego symulatora.
\tofill{Jedno zdanie z głównym wynikiem po zakończeniu badań.}

\section{Tło}

\paragraph{Zadanie.}
Iteracyjne przeszukiwanie lokalne~\cite{lourenco2003} dla permutacyjnego
problemu przepływowego na instancjach Taillarda~\cite{taillard1993} przechodzi
do najlepszej permutacji w wykładniczym sąsiedztwie: Adjacent,
Fibonacciego~\cite{bozejko2026new}, Dynasearch albo
Motzkina~\cite{wojewodzki2026soco}. Ruch jest rozwiązaniem
QUBO~\cite{lucas2014ising,glover2019}, którego zmienne wybierają zamiany; dla
dużych permutacji przeszukiwanie rozkłada problem na nakładające się okna
i rozwiązuje każde na wyżarzaczu~\cite{wojewodzki2026windowed}. Wcześniejsze
prace stosowały urządzenia D-Wave także do szeregowania
jednomaszynowego~\cite{bozejko2022distributed}.

\paragraph{Urządzenie.}
D-Wave Advantage (system 4.1, topologia Pegasus) programuje współczynniki
Isinga $h_i$ i $J_{ij}$ problemu na nadprzewodzących kubitach strumieniowych.
Zmienna logiczna, której sprzężenia nie mieszczą się w grafie sprzętu, staje
się łańcuchem fizycznych kubitów spiętych silnym sprzężeniem
ferromagnetycznym, czyli zanurzeniem minorowym~\cite{choi2008minor}; łańcuchy
mogą się zrywać i wtedy rozwiązanie wymaga naprawy. Zaprogramowane
współczynniki obarczone są zintegrowanymi błędami sterowania, czyli
gaussowskimi odchyleniami $h$ i $J$ o części systematycznej
i losowej~\cite{dwaveice}, a urządzenie wyżarza w skończonej temperaturze.
Nowsze eksperymenty pokazują koherentną dynamikę kwantową w programowalnych
łańcuchach Isinga~\cite{king2022coherent}.

\section{Symulatory}

\paragraph{Dokładny stan podstawowy.}
Dla okien do około dwudziestu zmiennych pełne przeszukanie daje optimum
każdego QUBO. Pokazuje, co zwróciłby idealny wyżarzacz, i służy jako
odniesienie dla pozostałych.

\paragraph{Symulowane wyżarzanie.}
Symulowane wyżarzanie~\cite{kirkpatrick1983sa} próbkuje rozkład Boltzmanna
przy malejącej temperaturze krokami Metropolisa. Klasyczne próbniki
D-Wave~\cite{dwavesamplers} implementują je, razem z przeszukiwaniem tabu, dla
tego samego formatu problemu co urządzenie, co czyni je naturalnym
klasycznym zamiennikiem.

\paragraph{Symulowane wyżarzanie kwantowe.}
Symulowane wyżarzanie kwantowe próbkuje model Isinga w polu poprzecznym metodą
Monte Carlo po całkach po trajektoriach: układ kwantowy odwzorowuje się na
klasyczny z dodatkowym wymiarem czasu urojonego złożonym z warstw Trottera,
a pole poprzeczne maleje w czasie przebiegu~\cite{santoro2002qa,martonak2002pimc}.
Metoda odtwarza statystykę równowagową wyżarzania kwantowego, a nie jego
dynamikę w czasie rzeczywistym, i była punktem odniesienia przy ocenie
wczesnych urządzeń D-Wave~\cite{boixo2014evidence}. OpenJij~\cite{openjij}
udostępnia zarówno symulowane wyżarzanie, jak i symulowane wyżarzanie
kwantowe z konfigurowalnymi harmonogramami.

\paragraph{Dokładna dynamika.}
QuTiP~\cite{johansson2013qutip,lambert2024qutip5} całkuje równanie
Schr\"{o}dingera albo równanie master hamiltonianu wyżarzania wzdłuż zadanego
harmonogramu, uwzględniając dekoherencję wprost. Koszt rośnie wykładniczo
z liczbą kubitów, więc stosujemy go tylko na kilku małych oknach.

\paragraph{Metody inspirowane kwantowo.}
Symulowana bifurkacja~\cite{goto2019sb} i cyfrowy wyżarzacz~\cite{aramon2019da}
rozwiązują QUBO dynamiką inspirowaną układami fizycznymi. Nie symulują
wyżarzacza kwantowego, a część z nich jest komercyjna, ale często go
zastępują; traktujemy je jako opcjonalne punkty odniesienia.

\begin{table}[t]
\caption{Symulatory wyżarzania. Dwie ostatnie kolumny pochodzą z naszych pomiarów.}
\label{tab:sims}
\centering
\small
\begin{tabular}{lllcc}
\toprule
Symulator & Co liczy & Dostępność & maks.\ zmiennych & s / okno \\
\midrule
Pełne przeszukanie & dokładny stan podstawowy & własny kod & \tofill{?} & \tofill{?} \\
dwave-samplers (SA, tabu) & wyżarzanie klasyczne & otwarty & \tofill{?} & \tofill{?} \\
OpenJij (SA, SQA) & wyżarzanie klasyczne i kwantowe Monte Carlo & otwarty & \tofill{?} & \tofill{?} \\
QuTiP & dokładna dynamika (układ otwarty) & otwarty & \tofill{?} & \tofill{?} \\
Bifurkacja, cyfrowy wyżarzacz & heurystyki inspirowane kwantowo & częściowo komercyjne & \tofill{?} & \tofill{?} \\
\bottomrule
\end{tabular}
\end{table}

\section{Metodyka}

\paragraph{Scenariusze.}
W A1 symulator zwraca dokładny stan podstawowy każdego okna. W A2 wykonuje
dynamikę wyżarzania bez błędów sterowania: symulowane wyżarzanie i symulowane
wyżarzanie kwantowe z harmonogramem dopasowanym do urządzenia (czas
wyżarzania, temperatura efektywna) oraz dokładną dynamikę na małych oknach. A3
dodaje błędy sterowania urządzenia jako gaussowskie zaburzenia $h$ i $J$
i rozwiązuje problem przez to samo zanurzenie minorowe co urządzenie, więc
łańcuchy mogą się zrywać.

\paragraph{Porównanie z urządzeniem.}
Przebiegi na urządzeniu nie zapisały pełnych zbiorów próbek, więc porównanie
działa na poziomie całych przebiegów: każdy symulator zastępuje urządzenie
wewnątrz iteracyjnego przeszukiwania lokalnego na tych samych instancjach
i oknach, z tą samą liczbą wywołań wyżarzacza co na urządzeniu, a względne
odchylenie długości uszeregowania od dolnego ograniczenia Taillarda
porównujemy testem znaków par i $\tau$ Kendalla po instancjach. Jeśli
pozwoli czas wyżarzacza, mały bieg z zapisem zbiorów próbek wybranych okien
da porównanie na poziomie rozkładów \tofill{?}.

\section{Plan eksperymentów}

Badanie przebiega w trzech etapach: każdy symulator na jednym oknie;
wszystkie na instancjach z $n = 20$; najwierniejsze na pełnym zbiorze
($n = 20$ do $500$), na który stać symulowane wyżarzanie i symulowane
wyżarzanie kwantowe.

\begin{table}[t]
\caption{Dane z D-Wave użyte jako odniesienie (D-Wave Advantage, system 4.1, Pegasus).}
\label{tab:hw}
\centering
\small
\begin{tabular}{ll}
\toprule
Pozycja & Wartość \\
\midrule
Przebiegi (ILS / SA) & 5650 (3250 / 2400), maj--sierpień 2026 \\
Instancje & $n = 20$: 5200 przebiegów; $n = 50$: 200; $n = 100$: 100; $n = 200$: 100; $n = 500$: 50 \\
Sąsiedztwa & Adjacent, Fibonacci, Dynasearch, Motzkin; warianty podstawowe i rozszerzone~\cite{wojewodzki2026windowed} \\
Wywołania wyżarzacza & 24\,317 (11 maja -- 8 sierpnia 2026) \\
Czas dostępu do QPU & łącznie 753 s \\
Rozmiar problemu na wywołanie & mediana 16 zmiennych, maksimum 220 \\
Zapisane dla przebiegu & historia długości uszeregowania, najlepsza permutacja, statystyki wywołań (bez zbiorów próbek) \\
\bottomrule
\end{tabular}
\end{table}

\section{Wyniki}

\tofill{Szkic: powstanie po trzech etapach.}

\begin{figure}[t]
\centering
\fbox{\parbox[c][3.5cm][c]{0.9\textwidth}{\centering \tofill{Figura: względne
odchylenie długości uszeregowania, symulator wobec D-Wave, jeden punkt na
instancję, dla A1, A2 (SA, SQA) i A3.}}}
\caption{Zgodność z urządzeniem na poziomie przebiegów.}
\label{fig:prd}
\end{figure}

\begin{table}[t]
\caption{Mediana względnego odchylenia [\%] przy tej samej liczbie wywołań wyżarzacza co na urządzeniu ($n = 20$).}
\label{tab:prd}
\centering
\small
\begin{tabular}{lccccc}
\toprule
Sąsiedztwo & D-Wave & A1 dokładnie & A2 SA & A2 SQA & A3 SQA + błędy \\
\midrule
Adjacent & \tofill{?} & \tofill{?} & \tofill{?} & \tofill{?} & \tofill{?} \\
Fibonacci & \tofill{?} & \tofill{?} & \tofill{?} & \tofill{?} & \tofill{?} \\
Dynasearch & \tofill{?} & \tofill{?} & \tofill{?} & \tofill{?} & \tofill{?} \\
Motzkin & \tofill{?} & \tofill{?} & \tofill{?} & \tofill{?} & \tofill{?} \\
\bottomrule
\end{tabular}
\end{table}

\section{Dyskusja}
\tofill{Szkic: który symulator może zastąpić urządzenie i do czego;
ograniczenia: brak zbiorów próbek z urządzenia, model błędów z dokumentacji,
jedna generacja urządzenia.}

\section{Wnioski}
\tofill{Szkic.}

\bibliographystyle{plain}
\bibliography{references}
\end{document}
